\documentclass[a4paper,11pt]{article}

\usepackage{luatexja}
\usepackage{luatexja-fontspec}
\setmainjfont{Noto Serif CJK JP}
\setsansjfont{Noto Sans CJK JP}
\setmonojfont{Noto Sans Mono CJK JP}
\usepackage[margin=20mm]{geometry}
\usepackage{amsmath}
\usepackage{booktabs}
\usepackage{array}
\usepackage{tabularx}
\usepackage{longtable}
\usepackage{enumitem}
\usepackage{graphicx}
\usepackage{xcolor}
\usepackage[hidelinks,unicode]{hyperref}
\emergencystretch=2em

\definecolor{statusgreen}{RGB}{32,115,70}
\definecolor{statusamber}{RGB}{167,102,0}
\definecolor{statusred}{RGB}{170,45,45}
\newcommand{\done}{\textcolor{statusgreen}{完了}}
\newcommand{\pending}{\textcolor{statusamber}{要確認}}
\newcommand{\todo}{\textcolor{statusred}{未完了}}
\newcommand{\code}[1]{\texttt{#1}}
\setlist{nosep,leftmargin=24pt}
\renewcommand{\arraystretch}{1.18}

\hypersetup{
  pdftitle={AJP simulation 総合現状・計算手法・移行計画},
  pdfauthor={AJP simulation workspace}
}

\title{AJP simulation\\総合現状・計算手法・移行計画}
\author{AJP simulation workspace}
\date{2026年9月9日}

\begin{document}
\maketitle
\tableofcontents
\clearpage

\section{この文書の位置づけ}

本書は、分散していた開発メモ、運用報告、障害調査、メッシュ検討、および有限半径
壁面沈着の導入手順を整理し、AJPノズル最適化計算の\textbf{現行正本}としてまとめた
ものである。以下を一冊で判断できることを目的とする。

\begin{itemize}
  \item ここまで何を実装・検証してきたか
  \item 現在のコードがどの計算手法を採用しているか
  \item どこまでを計算結果として信用でき、何が未検証か
  \item 次の本番キャンペーンより前に何をしなければならないか
\end{itemize}

CFD発散のログ分析、失敗点統計、同一メッシュA/B試験の全詳細は、別冊
\nolinkurl{cfd_failure_diagnosis_ja.tex}を根拠資料として残す。本書に記載するクラスタ上の
観測件数は過去レポートのスナップショットであり、XEON各ホストのライブ状態は
2026年9月9日時点で未確認である。

\section{エグゼクティブサマリー}

本プロジェクトは、11設計変数からAJPノズル形状・流量を生成し、Gmsh/OpenFOAMの
定常RANS、Lagrangian粒子追跡、3目的2制約の多目的ベイズ最適化を、4ホスト最大16枠で
非同期実行する段階まで拡張されている。

2026年9月8日には、従来不足していた有限粒径の壁面直接さえぎりを
\code{finiteRadiusDeposition}として追加した。OpenFOAM v2406での共有ライブラリbuild、
実行時load、強制接触1000 parcelのイベント記録・除去、unit testの通過までは
ローカルで確認した。したがって、\textbf{現行テンプレートの計算には固定壁面に対する
有限半径の直接さえぎり判定が入っている}。ただし、4 particle shard実行、各クラスタ
環境、時間刻み・メッシュ感度、および旧点粒子結果とのA/B比較は未完了である。

有限半径判定は粒子の運命と3目的値を変えるため、導入前の1,209件以上の履歴を新結果へ
混ぜてはならない。現行設定は
\code{central\_target\_d10um\_mobo\_v2\_structured\_finite\_radius}、campaignは
\code{mobo\_target10um\_v2\_structured\_finite\_radius}へ分離し、baselineを0にした。
旧設定と生成済み検証caseはarchiveへ退避した。次の本番計算前には、代表点を有限半径法で
再計算して新しいseedを作る必要がある。

production meshはstructured wedgeへ一本化した。正本は\code{base/meshGen2.py}であり、
通常ケース、BOケース、mesh sweepのすべてが同じgeneratorを使う。旧unstructured版は
\nolinkurl{archive/mesh/meshGen2_unstructured.py}へ移し、productionから参照しない。

同日、\code{potentialFoam}の全件SIGFPEの原因を、入口速度境界が未初期化のまま
\code{-writep}へ進んでいたことと特定した。\code{-initialiseUBCs -writep}へ修正し、
structured実meshで\code{potentialFoam}と100反復の\code{simpleFoam}が正常終了することを
確認した。\code{simpleFoam}には保守的な\code{residualControl}も追加し、10,000反復は
収束しない場合の上限として残した。

\section{これまでの流れ}

\begin{longtable}{p{21mm}p{82mm}p{48mm}}
\toprule
時期 & 主な作業 & 現在の扱い \\
\midrule
\endhead
2026-06-02 & 初期CFD・粒子計算ファイルを作成。 & 基礎として継続。 \\
2026-07-03 & SLURM上のケース生成、CFD、粒子、後処理を一連のpipelineとして整備。 & 継続使用。 \\
2026-07中旬 & scalar BO、中央直径10\,$\mu$m target、ケース回収・再提案を実装。 & 旧キャンペーン。 \\
2026-07-22 & 異常CFD場が粒子solverを停止させる問題を調査。CFD完了guard、粒子field検査、parser、注入位置を修正。 & コードへ反映済み。 \\
2026-07-23 & 大形状でunstructured meshが肥大化する問題を分析。structured generatorとsweepを作成。 & sweep合格。 \\
2026-07-24 & XEON6/1/4/5への分散dispatcherを整備。各4枠、最大16件を非同期運用。 & 設計は現行。ライブ状態は要確認。 \\
2026-07-27 & scalar v1、lexicographic v2を経て、3目的2制約のqLogNEHVIへ移行。旧履歴1,209件を再処理。 & 有限半径導入前の履歴。 \\
2026-07-27 & CFD失敗を統計・ログ・同一メッシュA/Bで診断。圧力をequation緩和からfield緩和へ変更。 & 修正反映済み。 \\
2026-09-08 & 有限半径壁面沈着を実装し、直接さえぎりを追加。 & ローカルserial検証済み。本番移行前。 \\
2026-09-08 & meshをstructuredへ一本化。\code{potentialFoam}の入口BC初期化を修正し、\code{residualControl}を導入。 & 実mesh smoke test済み。 \\
2026-09-08 & 新規計算に向け、旧configと2.9\,GBの生成済み\code{test\_cases}をarchive。active campaign/versionを有限半径版へ分離しbaselineを0化。 & 新campaignは未投入。 \\
2026-09-09 & GitHub共有用に利用者固有設定を\code{site.env}へ分離。ユーザー名、worker領域、coordinator用Pythonを一か所で変更可能にした。 & 設定診断・回帰テストを追加。 \\
\bottomrule
\end{longtable}

この流れを整理すると、これまでの中心課題は「多数形状を止めずに評価する基盤の構築」
から、「評価値を物理的・数値的に一貫させて新しいMOBOを開始すること」へ移っている。

\section{現行計算フローと正本コード}

\begin{center}
\begin{tabularx}{\textwidth}{>{\bfseries}p{34mm}X}
\toprule
処理 & 正本となる実装 \\
\midrule
候補生成・目的評価 & \code{bo\_loop.py}, \code{mobo.py}, \code{bo\_config.json} \\
4ホスト制御 & \code{tools/bo\_multihost.py}, \code{bo\_multihost\_config.json} \\
CFD template & \code{base/}, 特に\code{base/run.sh}と\code{base/system/} \\
メッシュ & structured正本: \code{base/meshGen2.py}; 旧unstructured: \code{archive/mesh/meshGen2\_unstructured.py} \\
粒子template & \code{baseparticle/}, 特に\code{loopaerosolDynamics.sh} \\
有限半径接触 & \code{baseparticle/custom/finiteRadiusDeposition/} \\
粒子集計 & \code{particle\_fates\_all.csv}生成部と\code{bo\_loop.py} \\
\bottomrule
\end{tabularx}
\end{center}

処理の流れは次の通りである。

\begin{enumerate}
  \item 11次元の候補を生成し、派生寸法・流量と幾何制約を計算する。
  \item Gmshでstructured 5度wedge meshを作成し、OpenFOAMへ変換する。
  \item \code{potentialFoam -initialiseUBCs -writep}で初期場を作る。失敗時はcaseを停止する。
  \item \code{simpleFoam}で定常RANSを解き、残差収束または10,000反復で停止する。
  \item CFD終了印、ログの\code{End}、最大速度を検査し、異常場を粒子段へ渡さない。
  \item \code{aerosolDynamicsFoam}で1000 parcelを追跡する。
  \item 標準の中心軌道壁面衝突と、有限半径の直接さえぎりを集計する。
  \item outlet、substrate、各nozzle wall、未解決へ分類する。
  \item 3目的・2制約へ変換し、制約付きqLogNEHVIで次候補を提案する。
  \item coordinatorが空いたホスト枠へ投入し、結果を回収する。
\end{enumerate}

\section{設計空間と幾何制約}

直接操作する変数と現在の範囲を表\ref{tab:variables}に示す。派生寸法は
\code{bo\_loop.py}内の変換関数で計算される。

\begin{table}[htbp]
\centering
\caption{11設計変数}
\label{tab:variables}
\begin{tabular}{lll}
\toprule
変数 & 範囲 & 概要 \\
\midrule
\code{R\_mm} & 0.5--2.0 & 基準半径 \\
\code{alphaRin} & 0.3--0.7 & 内側入口半径比 \\
\code{alphaL0} & 1.0--5.0 & 上流長さ比 \\
\code{L1\_mm} & 1.0--20.0 & 軸方向長さ \\
\code{alphaL2} & 0.5--5.0 & 下流長さ比 \\
\code{alphaL3} & 0.0--1.0 & 追加長さ比 \\
\code{theta0\_deg} & 10--60 deg & 第1傾斜角 \\
\code{theta1\_deg} & 10--60 deg & 第2傾斜角 \\
\code{alphaTheta2} & 0.0--1.0 & 第3角度比 \\
\code{alphat} & 0.5--5.0 & 厚さ比 \\
\code{U\_mps} & 1--50 m/s & 基準入口速度 \\
\bottomrule
\end{tabular}
\end{table}

生成後には、軸方向全長$\leq120$ mm、$L_2\leq70$ mm、$l_2\leq45$ mm、
$l_{2y}\leq45$ mm、$R_3\leq45$ mmを課す。これらは物理的実現性だけでなく、極端に
大きいmeshを候補生成前に抑える計算費用上の制約でもある。

\section{メッシュ生成}

\subsection{共通仕様}

軸対称断面を$y$軸まわりに5度回転した一層wedgeを用いる。境界patchは
\code{inletAerosol}, \code{inletSheath}, \code{outlet},
\code{wallUpper}, \code{wallDown}, \code{wallCavity},
\code{wallSubstrate}, \code{front}, \code{back}である。production環境変数でmeshを
120,000 nodes、360,000 elements以下に制限している。

\subsection{structuredへの一本化}

structured版は断面をblock状に分割し、軸方向・半径方向の分割数を制御する。旧
unstructured版は距離field細分化によって大形状でcell数が急増したため、productionから
外してarchiveへ移した。

structured sweepではedge 12/12形状がwedge角1度と5度の両方で
\code{checkMesh}に合格し、5度時のcell数は19,635--77,537、最大non-orthogonalityは
59.6808、最大skewnessは1.52498であった。追加random 10/10形状も
\code{checkMesh}と短時間\code{simpleFoam}に合格した。2026年9月8日以降は
\code{base/meshGen2.py}そのものをstructured正本とし、BO側の条件付き上書き処理も削除した。
今後のmethod metadataには\code{mesh=structured}とgenerator revisionを常に記録する。

\section{搬送流CFD}

\subsection{物理モデルと境界条件}

carrierは非圧縮・定常の\code{simpleFoam}で解き、RAS乱流モデルに\code{kOmegaSST}を
使う。入口のaerosol流とsheath流は\code{flowRateInletVelocity}で与え、5度wedgeに対応する
流量へ換算する。固体壁は\code{noSlip}、front/backは\code{wedge}とする。

現行の主な数値設定は次である。

\begin{itemize}
  \item 圧力: \code{relaxationFactors/fields/p = 0.3}
  \item $U,k,\omega$: equation under-relaxation 0.7
  \item 残差停止: $p<10^{-5}$、$U<10^{-4}$、$k,\omega<10^{-5}$
  \item 最大反復: 10,000（残差条件を満たさない場合の上限）
  \item 初期乱流量: $k=10^{-4}$、$\omega=1$
\end{itemize}

\subsection{完了判定と既知問題}

\code{simpleFoam}の終了code、ログ末尾の\code{End}、
$\max|U|\leq10{,}000$ m/sを確認し、\code{CFD\_DONE}または\code{CFD\_FAILED}を置く。
このguardにより、数値発散した速度場を粒子solverへ渡さない。

2026年7月の失敗35件では、圧力をequation側で緩和した旧設定が直接原因であった。同一
meshの代表3件をfield緩和へ変更すると800反復を完走した。高い
\code{theta0\_deg}は不安定性を顕在化させたが、それ自体を不可能形状の根拠にはできない。
詳細は別冊\code{cfd\_failure\_diagnosis\_ja.pdf}を参照する。

\code{potentialFoam}は過去の調査対象179/179件でSIGFPEとなっていた。原因は、生成直後の
\code{U}で\code{flowRateInletVelocity}の保存値がゼロのままなのに、境界条件を評価せず
\code{-writep}へ進んだことである。全域$U=0$となり、圧力近似に使う流向filter
\[
 \boldsymbol{F}=\frac{\boldsymbol{U}\boldsymbol{U}}
 {|\boldsymbol{U}|^2+\mathrm{SMALL}\,\langle|\boldsymbol{U}|^2\rangle}
\]
の分母もゼロになってSIGFPEを起こしていた。

現行は\code{potentialFoam -initialiseUBCs -writep}とし、入口流量境界を評価してから解く。
structured実meshのlocal testでは、Phiと$p$を非ゼロ反復で解き、SIGFPEなしで\code{End}へ
到達した。続く\code{simpleFoam}も100反復を正常完走した。今後はpotential初期化の失敗を
無視せず、\code{CFD\_FAILED}として停止する。

\code{residualControl}は上記の保守的閾値で導入済みである。閾値を満たすケースだけを早期
終了し、満たさないケースは従来通り最大10,000反復まで進む。粒子目的値に対する停止閾値
感度はsmall batchで引き続き確認する。

\section{粒子追跡}

\subsection{基礎モデル}

vendored \code{aerosolDynamicsFoam}の\code{basicKinematicCloud}を用いるone-way uncoupled
Lagrangian trackingである。粒子速度$\boldsymbol{u}_p$はcarrier速度を
\code{cellPoint}補間し、現在の力モデルは\code{sphereDrag}のみである。

粒径は0.1, 0.2, 0.3, 0.5, 0.8, 1, 2, 3, 5, 7\,$\mu$mの10 bin、各bin 100位置、
合計1000 parcelである。\code{analyticPositionTracking true}、時間刻み上限$10^{-4}$ s、
集計chunk 0.1 sとする。95\%のfateが確定するか、
\[
 T_{\max}=\min\!\left(5\ \mathrm{s},\ \max\!\left(0.5\ \mathrm{s},\
 3T_{\mathrm{transit}}\right)\right)
\]
へ到達した時点で停止する。

4並列は通常の\code{decomposePar}ではない。各rankがfull meshを持ち、注入粒子を4 shardへ
分けて\code{-particleParallel -nParticleShards 4}で追跡し、後でfateを結合する。

\subsection{有限半径による直接さえぎり}

標準\code{standardWallInteraction/stick}は、粒子中心軌道がpatchへ到達した場合を捕捉する。
これに加え、共有ライブラリ\code{finiteRadiusDeposition}が各移動区間
$\boldsymbol{x}_0\rightarrow\boldsymbol{x}_1$に対して有効半径
\[
 r_{\mathrm{eff}}=f_r d_p+\delta
\]
を用いる。現行値は$f_r=0.5$、$\delta=0$であり、球形粒子の物理半径$d_p/2$に相当する。

アルゴリズムは次の通りである。

\begin{enumerate}
  \item 対象wall faceからoctreeを構築する。
  \item 移動segmentのbounding boxを$r_{\mathrm{eff}}$だけ広げ、候補faceを検索する。
  \item segmentとfaceの最小距離が$r_{\mathrm{eff}}$以下となる候補を抽出する。
  \item 各候補の最小距離位置までを32回二分探索し、最初に
    $\mathrm{dist}(\boldsymbol{x},\mathrm{wall})\leq r_{\mathrm{eff}}$となる位置を求める。
  \item 最も早い接触をpatch名、粒径、particle ID、中心・壁面座標とともに記録し、
    parcelを除去する。
\end{enumerate}

対象patchは\code{wallSubstrate}, \code{wallUpper}, \code{wallDown},
\code{wallCavity}、actionは\code{remove}である。イベントは各shardから
\code{particle\_fates\_all.csv}へ統合される。

\subsection{「さえぎり」に含むものと含まないもの}

本実装が含むのは、既に抗力で計算された粒子中心軌道に対し、固定壁面と有限半径球の
幾何学的初回接触を加える\textbf{直接さえぎり}である。中心が壁面を横切らなくても表面が
触れれば沈着する。

次は含まない。

\begin{itemize}
  \item 先に沈着した粒子が後続粒子を遮る効果
  \item 堆積層成長による壁面形状・carrier流れ場の変化
  \item 粒子同士の衝突・凝集
  \item 粒子からcarrierへの二方向連成
  \item Brownian/random forceおよび乱流分散
\end{itemize}

したがって、「さえぎりを入れた」ことと、堆積成長・目詰まりまでモデル化したことは
区別しなければならない。

\section{粒子fateと目的関数}

粒子は\code{outlet}, \code{wallSubstrate}, \code{wallUpper},
\code{wallDown}, \code{wallCavity}, unresolvedへ分類する。注入総数を$N_{\rm inj}$、
基板中心$(x,z)=(0,0)$の半径5\,$\mu$m円内を$N_{\rm target}$、基板到達総数を
$N_{\rm substrate}$、基板以外の3 wallへの沈着を$N_{\rm wall}$とする。現行の3目的は
すべて最大化である。

\begin{align}
 f_1 &= \frac{N_{\rm target}}{N_{\rm inj}},\\
 f_2 &=
 \begin{cases}
 N_{\rm target}/N_{\rm substrate}, & N_{\rm substrate}>0,\\
 0, & N_{\rm substrate}=0,
 \end{cases}\\
 f_3 &= \frac{5}{5+r_{90}[\mu\mathrm{m}]}
        \min\!\left(1,\frac{N_{\rm target}}{20}\right).
\end{align}

$r_{90}$は基板粒子位置全体についてtarget中心から測った距離の90\%分位である。第3目的は
狭い分布を好むだけでなく、target粒子が少ない場合の見かけ上の高compactnessをsupport項で
抑える。

制約は
\[
 \frac{N_{\rm resolved}}{N_{\rm inj}}\geq0.95,
 \qquad
 \frac{N_{\rm wall}}{N_{\rm inj}}\leq0.05
\]
である。有限半径イベントも同じfate CSVへ入るため、直接さえぎりの追加は特に
$N_{\rm wall}$、$N_{\rm substrate}$、$N_{\rm target}$を変え得る。

\section{多目的ベイズ最適化と分散実行}

現行はBoTorchの制約付き\code{qLogNEHVI}で、3目的のPareto frontとhypervolumeを扱う。
初期点が32件に満たない間は空間充填候補を用いる。主な設定は次である。

\begin{center}
\begin{tabular}{lr}
\toprule
項目 & 現行値 \\
\midrule
候補pool & 2048 \\
GP学習点上限 & 384 \\
recent学習点 & 64 \\
Monte Carlo sample & 64 \\
acquisition batch処理 & 128 \\
初期batch & 32 \\
1回の追加batch & 8 \\
最大pending & 16 \\
\bottomrule
\end{tabular}
\end{center}

XEON6がcoordinatorとなり、XEON6、XEON1、XEON4、XEON5の各4枠へ独立Slurm jobを
割り当てる。現行の新campaignはbaseline 0、新規最大256観測、最小96観測後に
hypervolume改善が96観測停滞すれば終了する。旧campaignのbaseline 1,209観測と設定は
\nolinkurl{archive/configs/}へ保存し、新GPへ読み込まない。

新campaign directoryにはgeneration 0のmaximin初期候補32件を作成済みである。全候補が
幾何・計算量制約を満たし、観測0、Slurm投入0、host割当0の準備状態にある。観測数が
\code{gp\_min\_points=32}未満の間はPyTorch/BoTorchをloadせず、空間充填maximinを明示的に
使う。32件の評価後に初めてqLogNEHVIへ移る。

各hostでOpenFOAM版が同一とは限らない。共有ライブラリのbinaryを使い回さず、同じsourceと
辞書設定から、各hostの実行solverと同じOpenFOAM ABIで個別buildする必要がある。

\subsection{GitHub共有と利用者ごとの設定}

Linux login名やhome directoryを共有configへ直書きしない。GitHubで追跡する
\code{site.env.example}を各利用者が\code{site.env}へコピーし、次の利用者固有値を初回に設定する。

\begin{itemize}
  \item \code{AJP\_MASTER\_PYTHON}: XEON6 coordinator用Python
  \item \code{AJP\_SSH\_USER\_XEON1/4/5}: remote hostごとのSSH login名。host間で異なる値を許す。
  \item \code{AJP\_WORKER\_ROOT\_XEON6/1/4/5}: hostごとの利用者用worker領域の絶対path
\end{itemize}

\code{site.env}は\code{.gitignore}の対象であり、commitしない。共有する計算設定
\nolinkurl{bo_config.json}と\nolinkurl{bo_multihost_config.json}には、物理条件、campaign、
host addressを残す。dispatcherはhost別の環境変数を展開し、SSH targetを各hostの利用者名と
addressから組み立てる。remote case pathとSlurmへ渡す\code{AJP\_WORKER\_ROOT}もhostごとに
切り替える。この分離により、学生は共有configを書き換えずに同じ計算を実行でき、個人pathの
誤commitも避けられる。

各host固有のOpenFOAM環境は\code{tools/worker-envs/<host>.sh}を共有正本とし、実行先の
各hostの\code{AJP\_WORKER\_ROOT\_<HOST>/worker-env.sh}へ配置する。template中の利用者pathは
\code{\$HOME}から決まる。一方、OpenFOAMのinstall pathとABIはhost固有なので、cluster更新時は
template確認とcustom libraryの再buildが必要である。

\section{問題、原因、対策の整理}

\begin{longtable}{p{40mm}p{59mm}p{52mm}}
\toprule
問題 & 判明した原因・意味 & 対策と状態 \\
\midrule
\endhead
CFD発散後に粒子計算が止まる & 異常$U$場を粒子solverへ渡していた。 & 完了guardとfield検査を追加。\done \\
代表CFDがSIGFPE/異常速度 & 圧力equation under-relaxationで連続の式誤差が先行増大。 & \code{fields/p=0.3}へ変更し、代表3件で安定化。\done \\
$k,\omega$初期値がほぼゼロ & 乱流モデルの数値条件を悪化。主発散原因とは別。 & $k=10^{-4}$、$\omega=1$へ変更。\done \\
大形状でmesh肥大 & unstructured細分化が形状scaleに追随してcell数増大。 & structuredへ一本化し、旧版をarchive。\done \\
\code{potentialFoam}が全件SIGFPE & 未初期化の入口BCにより全域$U=0$のまま\code{-writep}の流向filterを評価。 & \code{-initialiseUBCs}を追加し、失敗をfatal化。実meshで検証済み。\done \\
\code{simpleFoam}が常に10,000反復 & 収束済みでもendTimeまで継続。 & 保守的\code{residualControl}を追加。目的値感度は\pending \\
粒子がcell境界付近で停滞 & wedge drift、境界chattering、短すぎるtracking horizon。 & analytic tracking、補正、設計別horizon。\done \\
点粒子判定が有限粒径を無視 & 中心軌道だけでは壁近傍を通る粒子の直接さえぎりを落とす。 & finite-radius CFOを実装。cluster/shard検証は\pending \\
報告書が時系列不明・重複 & 調査ごとのMarkdown、TeX、ログが並立。 & 本書へ統合し、根拠別冊以外を整理。\done \\
\bottomrule
\end{longtable}

\section{現状の信頼度}

\begin{center}
\begin{tabularx}{\textwidth}{p{43mm}p{24mm}X}
\toprule
項目 & 状態 & 判断 \\
\midrule
pipelineの構成 & \done & CFD、粒子、集計、MOBO、分散dispatcherがコード化済み。 \\
CFD発散guard & \done & 異常場を粒子段へ渡さない。 \\
圧力緩和修正 & \done & 同一mesh 3ケースのA/Bで有効。 \\
structured mesh production & \done & 正本を一本化し、edge 12/12、random 10/10、実solver smokeに合格。 \\
\code{potentialFoam}修正 & \done & structured実meshでPhi/$p$を解き、SIGFPEなしで終了。 \\
\code{residualControl} & \pending & 実solverで構文・100反復を確認。粒子目的値の閾値感度は未確認。 \\
有限半径CFO単体 & \done & v2406 build/load、強制1000接触、unit testまで確認。 \\
有限半径4 shard & \todo & ローカル環境ではMPI/PMIx制約により未確認。 \\
有限半径全host & \todo & hostごとのABI buildとsmall batchが必要。 \\
旧観測との分離 & \done & 新version/campaign、別directory、baseline 0へ変更。 \\
有限半径版の新seed & \todo & 旧代表点を新手法で再計算する必要がある。 \\
ライブcampaign状態 & \todo & coordinator正本のstatus取得が必要。 \\
\bottomrule
\end{tabularx}
\end{center}

\section{次に必要な作業}

\subsection{P0: 次の本番計算より前}

\begin{enumerate}
  \item \textbf{旧campaignをfreezeする。} XEON6のcoordinatorを正本として、実行中・投入済み・
    未回収ケース、観測件数、停止理由を取得する。新候補が継続投入されている場合は、
    version方針が決まるまで提案を止める。
  \item \textbf{active versionの分離を維持する。} 現在は
    \nolinkurl{central_target_d10um_mobo_v2_structured_finite_radius}、baseline 0へ変更済みである。
    今後もmesh種別、finite-radius係数、粒子solver/source revisionをmetadataへ記録する。
  \item \textbf{各hostでcustom libraryをbuildする。} binary copyではなく、各hostの
    solverと同じOpenFOAM環境で\code{baseparticle/build\_custom.sh}を実行する。
  \item \textbf{有限半径の基準検証を行う。} 同一caseで\code{radiusFactor=0}対0.5、
    serial対4 shard、時間刻み、meshを比較する。粒子収支、重複eventなし、
    $|\mathrm{center-wall}|\simeq d/2$を合格条件とする。
  \item \textbf{新seedを作る。} 旧Pareto点、代表点、wall沈着が多い点、粒径感度が高い点を
    有限半径法で再計算する。旧1,209件は比較用履歴として保持し、新GPへ直接混ぜない。
  \item \textbf{structured meshをversionへ記録する。} 選択肢には戻さず、generator revisionを
    campaign metadataへ固定する。
\end{enumerate}

\subsection{P1: 信頼性と計算費用}

\begin{enumerate}
  \item structured代表caseで流量、圧力損失、$U$場、粒子fateの再現性を確認する。
  \item residual停止点と3000/5000/10000反復の粒子評価を比較し、現在の
    $10^{-5}/10^{-4}/10^{-5}$閾値を確定する。
  \item residual条件とは別に、continuity errorや$|U|$の発散watchdogを追加する。
  \item \code{checkMesh -allGeometry -allTopology}のproduction gateを定義する。
  \item small finite-radius MOBO batchを動かし、dispatcher回収、fate集計、objective version
    filter、停止条件をend-to-endで検証する。
\end{enumerate}

\subsection{P2: 物理モデルの拡張}

\begin{itemize}
  \item 0.1--7\,$\mu$mの範囲でBrownian/random forceおよび乱流分散の寄与を比較する。
  \item 最終Pareto候補のmesh・粒子時間刻み独立性を確認する。
  \item 高フィデリティ計算または実験と比較し、低フィデリティmodelのbiasを定量化する。
  \item 必要性が確認された場合のみ、堆積層成長、粒子間相互作用、二方向連成を別versionで
    導入する。
\end{itemize}

有限半径とCFD停止条件を同じ検証batchで同時に変えると、結果差の原因を分離できない。
structured meshを固定したうえで変更を一つずつA/Bし、そのたびにmethod versionと設定を
保存する。

\section{再現・確認コマンド}

\begin{verbatim}
# 利用者ごとの初回設定（site.envはGit管理外）
cp site.env.example site.env
${EDITOR:-vi} site.env
python3 tools/check_site_config.py --config bo_multihost_config.json

# finite-radius library（実行solverと同じOpenFOAM環境で）
source /usr/lib/openfoam/openfoam2406/etc/bashrc
baseparticle/build_custom.sh

# unit tests
python3 -m unittest discover -s tests -v

# template設定
foamDictionary baseparticle/system/controlDict -entry libs -value
foamDictionary baseparticle/constant/kinematicCloudProperties \
  -entry cloudFunctions.finiteRadiusWall.radiusFactor -value

# coordinatorのライブ状態（XEON6で実行）
python tools/bo_multihost.py --config bo_multihost_config.json status
\end{verbatim}

有限半径validationでは、設定値だけでなく、生成された
\nolinkurl{finiteRadiusDeposition*.dat}と\nolinkurl{particle_fates_all.csv}の件数・ID・patchを
照合する。

\section{文書・証跡の管理方針}

今後のreport正本は本書\code{ajp\_project\_status\_ja.tex/pdf}とする。CFD障害の再検証に
必要な別冊\code{cfd\_failure\_diagnosis\_ja.tex/pdf}、その入力CSVと図だけを残す。新しい
調査を行う場合、日々のraw logを恒久reportとして増やさず、結論・条件・根拠fileを本書
または明確な別冊TeXへ反映する。

raw solver log、LaTeX中間生成物、同一内容のMarkdown/PDF、採用しない旧図は再生成可能な
作業物としてreport treeに残さない。数値条件を変更した際は、コードだけでなく本書の
更新日、method version、検証表を同時に更新する。

\code{tests/}は軽量なソース回帰テスト群であり、削除対象ではない。
\code{test\_cases/}はunit testではなく、mesh/solver検証で生成される大容量OpenFOAM caseの
出力先である。2026年7月分の26組、計2.9\,GBは
\nolinkurl{archive/test_cases_202607/}へ退避し、active側には用途説明だけを残した。
旧config 6件と旧移行script 2件も\nolinkurl{archive/configs/}および
\nolinkurl{archive/tools/}へ退避した。
さらに、現行multi-host BOから参照されない旧手動pipeline 10件と旧診断tool 4件を
\nolinkurl{archive/legacy_manual_pipeline_202607/}および
\nolinkurl{archive/tools/legacy_202607/}へ退避した。active rootにはBO本体・設定・templateへの
入口だけを置き、旧serial粒子runnerと旧mesh generatorはcaseへコピーしない。
利用者名と秘密鍵名を固定していた旧SSH configも同archiveへ退避し、active dispatcherは
Git管理外の\code{site.env}と共有host addressから接続先を構成する。

\section{当面の完了条件}

次のfinite-radius本番campaignを開始できる条件は、次の全項目である。

\begin{itemize}
  \item 旧campaignの投入がfreezeされ、観測snapshotが保存されている。
  \item 新method/objective versionが設定され、新campaign directoryが旧履歴と分離されている。
  \item 4 hostすべてで対応ABIのcustom libraryがbuild・loadできる。
  \item radius 0/0.5、serial/4 shard、粒子収支、接触距離のvalidationが合格する。
  \item structured generator revisionが固定され、少数caseでCFDと粒子の再現確認が完了する。
  \item residual停止と固定反復の目的値差が許容範囲内である。
  \item 旧代表点の再計算から新seedが作られ、旧点粒子観測と分離されている。
\end{itemize}

この条件を満たすまでは、現状を「有限半径機能は実装済み、production campaignへの移行は
未完了」と表現するのが正確である。

\end{document}
